\section{Technical Background}
\label{sec:technical-background}

\subsection{Convex Free-Space Covers for GCS--B\'ezier}
\label{sec:convex-free-space}

Obstacle avoidance makes motion planning difficult because the free
configuration space $\mathcal{C}_{\mathrm{free}}$ is generally non-convex.
The GCS--B\'ezier approach addresses this difficulty by representing the
part of $\mathcal{C}_{\mathrm{free}}$ relevant to a query with a finite
collection of collision-free convex regions,
\[
  \mathcal{Q}=\{Q_i\}_{i=1}^{N_R}, \qquad
  Q_i \subseteq \mathcal{C}_{\mathrm{free}}.
\]
The regions may overlap. They induce a graph with one vertex per region and an
edge for each feasible transition between adjacent regions. The start and goal
are connected to the regions that contain them or can be reached from them.

Graphs of Convex Sets (GCS) jointly optimize the discrete route through this
graph and the continuous trajectory variables associated with the selected
regions~\cite{MarcucciEtAl2023GCS,MarcucciEtAl2023MotionPlanning}. In the
GCS--B\'ezier formulation, a trajectory segment is represented by B\'ezier
control points constrained to its assigned region. By the convex-hull property
of B\'ezier curves, this containment certifies collision freedom for the whole
segment rather than only for sampled points.

Consequently, decomposition is not merely a preprocessing step. It determines
the number of graph vertices $N_R$, the number of feasible transitions $N_E$,
and the size of the resulting convex optimization model. A cover with fewer
regions is not necessarily easier to solve: large or highly overlapping regions
can induce a dense transition graph. This motivates evaluating decomposition
methods through their downstream GCS model and trajectory, not through region
count alone.

\subsection{Approximate Convex Decomposition}
\label{sec:acd}

Approximate Convex Decomposition (ACD) is a geometry-driven method for
decomposing a non-convex polygon~\cite{LienAmato2006}. It recursively locates
the most significant concavity, inserts a cut to resolve it, and repeats until
the residual concavity of every component is below a user-selected tolerance
$\tau$. The tolerance controls the granularity of the decomposition: a smaller
$\tau$ resolves more geometric detail and usually produces more regions, whereas
a larger $\tau$ avoids cuts associated with minor features.

For the present study, ACD provides a direct and inexpensive way to construct
the safe-region representation for a planar environment. Its principal
trade-off is that resolving local geometric features can produce many regions
and, therefore, many possible GCS transitions. This makes ACD a useful baseline
for testing whether a faster decomposition stage offsets the cost of a larger
downstream graph.

\subsection{Visibility Clique Cover}
\label{sec:vcc}

Visibility Clique Cover (VCC) constructs a convex cover from the global
visibility structure of the free space~\cite{DeitsTedrake2023VCC}. It first
samples collision-free configurations and connects pairs whose straight-line
segment remains collision free. Cliques in this visibility graph identify groups
of mutually visible samples. Each clique initializes an ellipsoid, which is
inflated into a collision-free convex polytope using an IRIS-style
region-inflation step~\cite{DeitsTedrake2014}. Sampling and inflation are
repeated until the desired coverage is reached.

Because it groups mutually visible samples before inflation, VCC can produce a
small number of large regions and substantially reduce the GCS graph for some
environments. Its cost, however, depends on the sampling budget, visibility
tests, clique-cover computation, and coverage target. Sparse or unrepresentative
samples can also make narrow passages difficult to cover or connect. The
comparison in this paper therefore examines the complementary trade-off between
ACD's inexpensive local decomposition and VCC's potentially more compact but
costlier visibility-based cover.
